\begin{helper}[Strict active-slot delta exclusion]
Let \(G,e_0,k,T\), the generating coloring, and the ETT be as in
Extension-base Part I.  Fix the last noninitial Lean index \(q\), with
\(q+1\) equal to the number of rounds, and assume its positive slot is a PE
whose exact chosen vertex is a fixed terminal support vertex in \(V(T)\).
Assume that every displayed tree at an index \(r<q\) is elementary under its
preceding displayed coloring.  Also assume the following lower-rung
maximum-defective order assertion: whenever \(a<r\), \(\sigma\) is a proper
exactly-one-uncolored coloring stable at the displayed tree for \(r\),
\(\delta\) is defective there, and \(v\) is a \(\delta\)-boundary vertex
having an earlier \(\delta\)-boundary vertex, then, if \(v\in V(T)\), its
first-appearance index is at most \(a+1\), and slot \(a.\mathsf{succ}\) is
PE or SE with exact chosen vertex \(v_a\), one has
\(v=v_a\) or \(v\prec v_a\) in that displayed tree.

Let \(s<t\) be two members of
\(L=\noderef{ExtensionBaseActiveIndices}\), with exact chosen vertices
\(v_s,v_t\in V(T)\), and suppose the exact delta field at slot
\(t.\mathsf{succ}\) is \(\delta_t\).  Then
\[
 \delta_t\notin
 \overline{\varphi_s}\bigl(T_{s+1}\bigr).
\]
All edges in this statement and proof are individual edge identities, so
parallel edges are never identified.
\end{helper}
\begin{proof}
Suppose that \(\delta_t\) is missing at a vertex \(u\) of the displayed tree
at \(s\), using \noderef{TreeMissingColors}.  Since \(s<t\le q\), the
elementarity hypothesis and \noderef{ColoredVertexSetElementary} imply that
every other vertex \(x\) of this tree has \(\delta_t\) present.  By
\noderef{MissingColors} and \noderef{PresentColors}, choose an exact edge
identity \(f_x\) incident with \(x\) and colored \(\delta_t\) by the
coloring at \(s\).

Membership of \(s\) in the active set forces its certificate to be PE.  Its
preceding-closed field, together with the missing witness \(u\), says that
\(f_x\) cannot be a boundary edge.  Since \(f_x\) is incident with the
inside vertex \(x\), the endpoint definitions \noderef{Incident} and
\noderef{BoundaryEdges} then put every endpoint of this same identity inside
the tree; hence \(f_x\) is in \noderef{InternalEdges}.  Apply
\noderef{SameColorSeriesPreserved} from \(s\) through the slot immediately
before \(t\).  Its internal-edge clause preserves the exact equation
\(\varphi_t(f_x)=\delta_t\).  Iterating both inclusions of
\noderef{HkuMissingStoredMonotone} over the same interval puts
\(\delta_t\) in the union of the missing and stored colors at the displayed
tree for \(t\).

Eliminate the PE certificate at \(t\).  Let \(\pi\) be its selected coloring
and \(\rho\) its second stable coloring.  The certified ordered extensions,
expressed by \noderef{TreeSequenceInitialSegment}, include the tree at
\(s\) in the tree at \(t\).  The edge clause of
\noderef{StableColoring}, first for \(\pi\) relative to \(\varphi_t\) and
then for \(\rho\) relative to \(\pi\) with \(\delta_t\) inserted in the
controlling set, gives
\[
             \pi(f_x)=\rho(f_x)=\delta_t.                 \tag{1}
\]
Consequently \(x\) is not a \(\delta_t\)-boundary vertex under \(\rho\).
Indeed, otherwise \noderef{BoundaryColorVertex} supplies an exact boundary
identity \(g\) incident with \(x\) and colored \(\delta_t\).  The internal
identity \(f_x\) is not a boundary identity, so \(g\ne f_x\).  These two
distinct identities are adjacent by \noderef{AdjacentEdges}, and (1)
contradicts the properness condition \noderef{IsProperEdgeColoring} for
\(\rho\).  Thus every vertex of the old tree other than \(u\) is excluded
from the \(\delta_t\)-boundary vertices under \(\rho\).                 \(\tag{2}\)

Restrict the enlarged stability of \(\rho\) to the old stored set.  Its
vertex clause is unchanged and its edge clause only loses controlling
colors.  Compose this restricted relation with the selected coloring's
stability by the transitivity clause of \noderef{sc2}.  Hence \(\rho\) is
proper, leaves exactly \(e_0\) uncolored, and is stable relative to the
preceding displayed coloring at slot \(t\), exactly as required by
\noderef{MaximumDefectiveChoice}.

The exact \(\delta_t\)-boundary-edge set is the same under \(\pi\) and
\(\rho\).  Membership in \noderef{BoundaryEdges} is coloring-independent.
For a boundary identity colored \(\delta_t\) under \(\pi\), the enlarged
edge-stability clause fixes it under \(\rho\).  Conversely, the symmetric
criterion \noderef{sc1} fixes a boundary identity colored \(\delta_t\)
under \(\rho\).  Thus \noderef{BoundaryColorEdges} gives equality of the
two finite sets.  Expanding \noderef{DefectiveEdge} for the selected exact
edge shows that its defective color is precisely the recorded
\(\delta_t\), by comparison with the stored equation on that edge.
Therefore \noderef{DefectiveColor} supplies a second, distinct exact
\(\delta_t\)-boundary identity \(g\) under \(\rho\).

Choose its inside endpoint \(y\).  Then \(y\) and the selected vertex
\(v_t\) are \(\delta_t\)-boundary vertices.  They are distinct, since if
they coincided, the distinct identities \(g\) and the selected edge would
be adjacent and have the same color, again contradicting properness.  The
vertex \(y\) is defective under \(\rho\).  Maximality in
\noderef{MaximumDefectiveChoice}, applied with the admissibility just
proved, yields
\[
                    y\prec_{T_{t+1}}v_t.                    \tag{3}
\]

Apply \noderef{ExtensionBaseSupportOrder} with the stated lower-rung order
hypothesis.  It gives \(v_t=v_s\) or \(v_t\prec_Tv_s\).  Iterating the
ordered-extension fields of \noderef{TashkinovSeriesStep} and the terminal
segment field of \noderef{ExtendedTashkinovTree} shows that both displayed
trees are ordered initial segments of \(T\).  The entrywise prefix equations
preserve \noderef{TreeVertexOrder}; moreover, an earlier vertex of \(T\)
than a vertex in an ordered prefix is itself in that prefix, by injectivity
of the vertex sequence in \noderef{TreeSequence}.  From (3) it follows that
both \(v_t\) and \(y\) lie in the displayed tree at \(s\).

If \(u=v_t\), then the boundary vertex \(y\) belongs to the old tree and is
different from \(u\), contradicting (2).  If \(u\ne v_t\), then the boundary
vertex \(v_t\) itself contradicts (2).  The assumed missing witness cannot
exist, proving the asserted strict-slot exclusion.
\end{proof}
